\subsection{Two further boundaries for changed scalar coefficients}

The preceding scalar-complex criterion gives two sharper
tests. One excludes a concrete nonflat family even after
changing the coefficient complex; the other excludes the
natural pure coefficient realization without a flatness
hypothesis on the unknown lift.

\begin{finding}[\fid{F59}, the singular invariant factor]
\label{f:scalarInvariantFactorNoGo}
Let \(k\) have characteristic zero and put
\[
 A=k[x,y,z_1,\ldots,z_m],\qquad
 \tau(x,y,z)=(-x,-y,z),\qquad
 S=A^\tau,\qquad T=A\otimes_S A.
\]
The variables have positive even degree, with \(x,y\)
of degree two. Let \(U\) be a bounded complex of finite
graded \(k\)-spaces with nonzero cohomology. For
\[
                         X=U\otimes_kT,
\]
there is no bounded scalar complex \(C\) and actual
derived bimodule isomorphism
\begin{equation}\label{eq:scalarInvariantChangedQuadratic}
                 X\otimes_A^{\mathbf L}X\simeq C\otimes_kX.
\end{equation}
This conclusion requires no prescribed Euler class.
\end{finding}

\begin{proof}
Generically \(T\) has one identity graph and one
\(\tau\)-graph, by the explicit invariant-correspondence
calculation above. Forget internal grading and set
\[
 u(t)=\sum_i\dim_kH^i(U)t^i,\qquad
 d(t)=\sum_i\dim_kH^i(C)t^i.
\]
Generic K\"unneth and the proposed quadratic give
\[
             2u(t)^2(e+\tau)=d(t)u(t)(e+\tau).
\]
Since \(u\ne0\), the Laurent polynomial domain gives
\(d=2u\). Let \(a\) be the highest cohomological degree
of \(U\), and let \(r=\dim_kH^a(U)>0\).
Then \(C\) has the same highest degree \(a\), with
\(\dim_kH^a(C)=2r\).

The ordinary module \(T\) is in degree zero. Its derived
square has top degree zero, with \(H^0=T\otimes_AT\).
That module is nonzero because it has generic rank four.
Taking top degree \(2a\) in
\eqref{eq:scalarInvariantChangedQuadratic} gives the
actual module isomorphism, after forgetting internal grading,
\begin{equation}\label{eq:scalarInvariantTopDegree}
             (T\otimes_AT)^{\oplus r^2}
                         \simeq T^{\oplus2r^2}.
\end{equation}
This uses only top degrees and assumes no vanishing of
the possible lower Tor groups.

At the origin of the left \(A\)-space,
\[
 B=k\otimes_AT
   =k[x,y]/(x^2,xy,y^2),\qquad \dim_kB=3.
\]
The fixed variables \(z_i\) are also killed. Every
positive invariant acts as zero on \(B\), so
\[
 \begin{aligned}
 k\otimes_A(T\otimes_AT)
   &=B\otimes_SA\\
   &=B\otimes_k(k\otimes_SA)
     =B\otimes_kB.
 \end{aligned}
\]
This has dimension nine. Applying the ordinary origin
fiber to the actual module isomorphism
\eqref{eq:scalarInvariantTopDegree} therefore gives
\(9r^2=6r^2\), impossible.
Ordinary fiber is being applied to an actual isomorphism,
not asserted to preserve an arbitrary derived triangle.
\end{proof}

This strengthens the fixed-coefficient origin-fiber test.
It does not assert that every coupled complex factors
as \(U\otimes_kT\). A genuinely nonfactorable nonflat
complex remains outside this finding.

\begin{finding}[\fid{F59}, the pure scalar coefficient boundary]
\label{f:pureScalarSoergelNoGo}
Retain the actual generator and scalar quadratic hypotheses
of Finding~\ref{f:alteredScalarSoergel}. For the arithmetic
clause take the scalar field \(k=E\), a finite
\(\ell\)-adic coefficient field with \(\ell\nmid q\),
enlarged as required.
Suppose the underlying scalar complex \(C\) has a Weil
structure over \(\mathbb F_q\), is bounded and pure of
one integer complex weight, and satisfies, for every
sufficiently large \(n\),
\begin{equation}\label{eq:pureScalarComparison}
 \operatorname{Tr}_{\mathrm{Eul}}(F^n;C)
  =\epsilon\,\beta^n(q^{sn}+2+q^{-sn}),
 \qquad s\in\mathbb Z\setminus\{0\},\quad
 \epsilon\in\{1,-1\}.
\end{equation}
Here \(\beta\) is a nonzero pure algebraic eigenvalue
of an integer weight \(t\). Then no such ordinary
bounded Soergel lift \(X\) exists, even if its cohomology
is nonflat.

The arithmetic comparison is additional data; it is not
deduced from the formal scalar Euler class.
Alternatively one may first make a specified scalar
extension preserving the dimensions and relation above.
No extension from an arbitrary characteristic-zero field
to an \(\ell\)-adic coefficient field is implicit.
\end{finding}

\begin{proof}
Let \(w\) be the pure complex weight. On the finite-field
point, BBD's criterion puts an eigenvalue of weight \(a\)
only in degree \(a-w\)
\cite[5.1.8, pp.~126--127]{BBD}.
Consequently the same eigenvalue cannot occur in
opposite parities and cancel in the Euler trace.
Jordan blocks do not change its trace multiplicity.

Finite power-sum uniqueness follows from the invertible
Vandermonde matrix on distinct nonzero roots. The same
argument applies to eventual equality.
Thus the roots in \eqref{eq:pureScalarComparison} have
Euler multiplicities \(\epsilon,2\epsilon,\epsilon\),
at
\[
                     \beta q^s,\qquad \beta,\qquad\beta q^{-s}.
\]
There can be no other root: purity places it in only
one degree, so a zero Euler multiplicity would force
zero ordinary multiplicity.

These three weights are \(t+2s,t,t-2s\). Their degrees
are
\[
                     t+2s-w,\qquad t-w,\qquad t-2s-w.
\]
They are distinct since \(s\ne0\), and their actual
cohomological dimensions are \(1,2,1\). The common
sign merely determines the parity of those degrees.
Finding~\ref{f:alteredScalarSoergel}, however, requires
every cohomological multiplicity of \(C\), after
forgetting internal grading, to be even.
The two outer multiplicities contradict that condition.
\end{proof}

\paragraph{The precise consequence for the pure kernel construction.}
Retain the original formal grading labels if the
\fid{F61} coefficient is proposed as a scalar complex
in the ordinary Hecke category. The two candidates are
\[
 \begin{aligned}
 C_{\mathrm{full}}
   &=E(-2)[-4]\oplus E^{\oplus2}\oplus E(2)[4],\\
 C_{\mathrm{half}}
   &=E(-1)[-2]\oplus E^{\oplus2}\oplus E(1)[2].
 \end{aligned}
\]
With those labels their Euler class is \(\delta^2\).
Their cohomological multiplicities are \(1,2,1\) in
degrees \(-4,0,4\) and \(-2,0,2\), respectively.
This already contradicts the required evenness.
Equivalently, the pure scalar finding applies with
\(s=2\) and \(s=1\), \(\beta=1\) and \(\epsilon=1\).
A common pure arithmetic line, Tate twist or
cohomological shift does not change the odd outer
multiplicities, whenever the changed candidate retains
the stated scalar Euler premise.

The same coefficients and their actual maps do work
in \fid{F61}'s different finite-state kernel category.
The present conclusion excludes this scalar-coefficient
transfer into a bounded ordinary Soergel lift of the
actual ambient generator class with the required
quadratic. It does not contradict \fid{F61} or exclude
every pure realization in another category.

There is also a grading-only version. If
\(H^j(C)_e\) is supported on a single affine degree
line \(j=ae+b\), with nonzero integral slope \(a\),
then its Euler class \(\delta^2\) forces the two
outer multiplicities to be one in distinct degrees.
There is no opposite-parity cancellation at a fixed
formal degree on that line. Such a coefficient also
violates the generic evenness condition.
This requires the displayed grading-to-degree comparison.

\paragraph{The remaining nonflat case.}
Neither finding supplies a general exclusion of every
nonfactorable mixed coefficient lift with nonflat
cohomology. Nonflat quasi-idempotents exist in simpler
settings: for a nonzero homogeneous polynomial \(f\)
of positive degree,
the central two-term Koszul complex \(K(f)\) satisfies
\[
 K(f)\otimes K(f)
   \simeq K(f)\otimes
             \bigl(k\oplus k(-\deg f)[1]\bigr).
\]
Replace the second Koszul generator by its difference
with the first to obtain this actual isomorphism.
Its generic image and Euler rank are both zero.
It therefore does not supply the nonzero generic
support \(2(e+\tau)\) required for the actual generator.

% Separate paragraph for Part III.
\paragraph{The two changed-coefficient tests in Grothendieck's language.}
The singular invariant correspondence and its origin
fibers use EGA~I, Theorem~1.3.7, p.~86,
Corollary~1.3.8, p.~87, and Theorem~1.4.1,
pp.~90--92, as the affine module/sheaf framework
\cite{EGA}. The top-degree fiber argument is proved
with its actual isomorphism hypothesis; no general
cohomology-and-base-change assertion replaces it.
Formal grading has the SGA~3, I, 4.7.2--3 action
framework and cote~90, scan~17. The effective-versus-virtual
coefficient distinction uses SGA~6, V, \S\S2--3,
with cote~133, scan~18 \cite{SGA3,SGA6,Archive}.
No \(\lambda\)-ring structure on the noncommutative
word algebra is asserted.

The scalar complexes lie in SGA~6, I,
Definition~4.7, pp.~121--122, and Example~4.8,
p.~122 \cite{SGA6}. BBD, 5.1.8, pp.~126--127,
is the exact point-complex purity convention
\cite{BBD}. The finite power-sum uniqueness and
integer multiplicity contradiction are proved here.
The additional arithmetic comparison is not supplied
by a source canonical lattice or an unpublished
archive statement.
